\documentclass{article}
\usepackage{algorithm, algorithmic, amsfonts, amsmath}

\begin{document}
\begin{algorithm}[!h]
\renewcommand{\algorithmicrequire}{\textbf{Input:}}
\renewcommand{\algorithmicensure}{\textbf{Output:}}
	\caption{PUAE (deep PU anomaly detection for contaminated unlabeled data)}
	\begin{algorithmic}[1]
	    \REQUIRE
	    Unlabeled set $\{x_{n}\}_{n=1}^{N}$ (contaminated with anomalies);
	    labeled anomaly set $\{\tilde{x}_{m}\}_{m=1}^{M}$;
	    autoencoder $E_{\theta},D_{\theta}$; anomaly ratio $\pi$
	    \ENSURE
	    Anomaly detector $\theta$
	    \STATE \textbf{Convention:} normal data are the negative class ($y=0$),
	    anomalies are the positive class ($y=1$) % the reverse of ordinary PU: what is labeled are the anomalies
	    \STATE Reconstruction error $\ell(x;\theta)=\|x-D_{\theta}(E_{\theta}(x))\|^{2}$
	    \STATE \textbf{Turn the error into a probability:}
	    \STATE \quad $p_{\theta}(y\mid x)=\begin{cases}
	    \exp(-\ell(x;\theta)) & y=0\\
	    1-\exp(-\ell(x;\theta)) & y=1\end{cases}$ % a small reconstruction error means a high probability of being normal
	    \STATE \quad $\ell_{BCE}(x,y;\theta)=(1-y)\,\ell(x;\theta)-y\log\!\left(1-e^{-\ell(x;\theta)}\right)$ % both terms are bounded below, so maximizing the reconstruction error cannot diverge
	    \STATE \textbf{Baseline (ABC)} assumes all unlabeled points are normal:
	    \STATE \quad $\mathcal{L}_{ABC}(\theta)=\dfrac{1}{N}\sum_{n}\ell_{BCE}(x_{n},0;\theta)
	    +\dfrac{1}{M}\sum_{m}\ell_{BCE}(\tilde{x}_{m},1;\theta)$
	    \STATE \textbf{PU correction:} the unlabeled set is contaminated, so rewrite the normal-data
	    term through the mixture identity and form an unbiased PU risk % the PU risk replaces the assumption that U is entirely normal
	    \WHILE{stopping criterion is not met}
	    \STATE Sample batches from the unlabeled and the labeled-anomaly sets
	    \STATE Update $\theta$ on the PU-corrected objective
	    \ENDWHILE
	    \STATE \emph{The same construction applies to Deep SVDD, giving PUSVDD} % the method is independent of the anomaly-detection backbone
		\RETURN $\theta$
	\end{algorithmic}
\end{algorithm}
% Takahashi, Iwata, Kumagai, Yamanaka, Yamashita, arXiv 2024, "Deep Positive-Unlabeled
% Anomaly Detection for Contaminated Unlabeled Data". Compiled from Section 3 and equations (3), (4) and (5).
% The same idea appears twice in this group's work, here and in the PU diffusion model of the following year (ICLR
% 2025): an unbounded quantity such as a reconstruction error or ELBO is first turned into a bounded BCE form
% (equation 4), and the PU risk correction is then applied. The key design matches that paper: maximizing the
% reconstruction error of anomalies directly would diverge, whereas -log(1 - exp(-l)) bounds both terms below and
% balances them without an extra hyperparameter. Backbone independence: the paper instantiates both an autoencoder
% (PUAE) and Deep SVDD (PUSVDD).
\end{document}
